\chapter{Đối xứng, biến quan hệ và liên kết biểu diễn}
\chapterlead{Một phép biểu diễn tốt không chỉ mô tả nghiệm; mô hình còn quyết định cách
bộ giải nhìn thấy các nghiệm tương đương và các xung đột cục bộ.
\emph{Symmetry breaking} (phá đối xứng) thu gọn không gian tìm kiếm, còn
\emph{channeling} (liên kết biểu diễn) nối các tiêu chí bổ trợ của cùng một
quyết định.}

\index{symmetry breaking}\index{lex-leader}
\index{relation variable}\index{order encoding}\index{channeling}

\flowdiagram{Quyết định gốc}{Nhóm đối xứng}{Đại diện chuẩn}
  {Biến quan hệ}{Mệnh đề cục bộ}

\section{Hai cấp đối xứng: mô hình và công thức}

Một đối xứng của bài toán tác động lên đối tượng miền: đổi tên khung chứa, máy,
nhóm, tuần, màu hoặc phản xạ hình học. Một đối xứng của \CNF là hoán vị biến
và literal bảo toàn tập mệnh đề. Hai nhóm này liên quan nhưng không trùng nhau.
Phép biểu diễn có thể:
\begin{itemize}
  \item làm một đối xứng miền biến mất vì không biểu diễn chỉ số thừa;
  \item làm một đối xứng miền trở nên khó phát hiện do biến phụ;
  \item tạo đối xứng phụ mới giữa các trạng thái tương đương;
  \item có đối xứng cú pháp không tương ứng với một phép biến đổi nghiệm dễ
        diễn giải.
\end{itemize}

Các vị từ phá đối xứng tổng quát xây một điều kiện chọn đại diện cho các quỹ
đạo của bài toán tìm kiếm \parencite{crawford1996symmetry}. Ở cấp \CNF, những
hệ như Shatter xây đồ thị tô màu của công thức, tìm tự đẳng cấu rồi sinh các
ràng buộc phá đối xứng \parencite{aloul2003shatter}. Ở cấp mô hình, người thiết
kế đã biết ngữ nghĩa nên có thể viết ràng buộc ngắn hơn, chẳng hạn tiền tố
khung chứa đã dùng. Hai cách bổ sung nhau: phá đối xứng miền được áp dụng trước;
phát hiện trên \CNF tìm phần đối xứng còn sót.

\begin{center}
\captionof{table}{Bốn cấp khai thác đối xứng.}
\label{tab:symmetry-levels}
\begin{tabularx}{\textwidth}{@{}>{\bfseries\sffamily}p{.2\textwidth}YYY@{}}
\toprule
Cấp & Đầu vào & Lợi thế & Hạn chế\\
\midrule
Mô hình & Đối tượng và phép đổi tên & Ràng buộc ngắn, có chứng minh miền &
Cần tri thức bài toán\\
\CNF tĩnh & Đồ thị công thức & Tự động, dùng cho mô hình bất kỳ &
Biến phụ có thể che đối xứng\\
Động & Trạng thái tìm kiếm & Khai thác đối xứng còn lại & Chi phí tích hợp bộ giải\\
Chuẩn hóa & Toàn quỹ đạo & Có thể giữ một đại diện & Ràng buộc lớn/lan truyền sâu\\
\bottomrule
\end{tabularx}
\end{center}

Các lớp ràng buộc phá đối xứng được ghép theo thứ tự. Ràng buộc miền làm thay
đổi nhóm tự đẳng cấu của đồ thị; bộ phá đối xứng tự động vì thế chạy trên công
thức đã có mệnh đề đơn vị và liên kết biểu diễn. Mỗi lớp bảo toàn ít nhất một
đại diện của các quỹ đạo còn lại, còn thí nghiệm phân tích thành phần (ablation study) đo tác động của từng
lớp lên hiệu năng
\parencite{gent2006symmetry}.

\section{Đối xứng như tác động lên nghiệm}

Gọi \(S\) là tập nghiệm khả thi và \(\Gamma\) là một nhóm các phép biến đổi
bảo toàn tính khả thi và giá trị mục tiêu. Hai nghiệm \(s,s'\in S\) tương
đương khi tồn tại \(\gamma\in\Gamma\) sao cho \(s'=\gamma(s)\). Tập
\[
  \operatorname{Orb}(s)=\set{\gamma(s)\suchthat\gamma\in\Gamma}
\]
là quỹ đạo của \(s\).

Phá đối xứng thêm công thức \(B\) sao cho mỗi quỹ đạo còn ít nhất một
đại diện:
\[
  \forall s\in S,\qquad
  \operatorname{Orb}(s)\cap\operatorname{Mod}(B)\ne\varnothing.
\]
Không nhất thiết phải giữ đúng một đại diện. Trong thực hành, một tập ràng buộc
phá đối xứng nhỏ, rẻ và phù hợp với mô hình thường hữu ích hơn việc cố gắng biểu diễn
một dạng chuẩn hoàn chỉnh \parencite{walsh2006symmetry,gent2006symmetry}.

\begin{figure}[tb]
\centering
\begin{tikzpicture}[satfig,node distance=17mm]
  \node[satstate,sattrue] (s1) {\(s_1\)};
  \node[satstate,right=of s1] (s2) {\(s_2\)};
  \node[satstate,below=of s2] (s3) {\(s_3\)};
  \node[satstate,left=of s3] (s4) {\(s_4\)};
  \draw[satdash] (s1)--node[above,satlabel]{đổi tên}(s2);
  \draw[satdash] (s2)--node[right,satlabel]{phản xạ}(s3);
  \draw[satdash] (s3)--node[below,satlabel]{đổi tên}(s4);
  \draw[satdash] (s4)--node[left,satlabel]{phản xạ}(s1);
  \node[draw=Teal,very thick,rounded corners=4pt,fit=(s1),inner sep=4pt,
        label={[satlabel,text=Teal]above left:đại diện chuẩn}] {};
  \node[draw=Rule,dashed,rounded corners=7pt,fit=(s1)(s2)(s3)(s4),
        inner sep=8pt,label={[satlabel]below right:\(\operatorname{Orb}(s_1)\)}] {};
\end{tikzpicture}
\caption{Một quỹ đạo gồm bốn cách biểu diễn tương đương. Ràng buộc phá đối
xứng hợp lệ giữ ít nhất một nút; phép chuẩn hóa mạnh hơn sẽ cố giữ đúng một nút.}
\label{fig:symmetry-orbit}
\end{figure}

\begin{designrule}{Chứng minh theo phép biến đổi}
Chứng minh một ràng buộc phá đối xứng bằng phép hoán vị hoặc phép phản xạ đưa mọi nghiệm bị
loại sang một nghiệm giữ lại. Ánh xạ này đồng thời xác định nhóm đối xứng và
đại diện được bảo toàn.
\end{designrule}

\section{Loại đối xứng ngay từ lựa chọn biến}

Biến chính có thể tạo hoặc loại đối xứng. Trong Social Golfer, biến
\(G_{i,r,t}\) chỉ ghi người chơi, nhóm và tuần. Nếu thêm chỉ số vị trí trong
nhóm, mỗi nhóm \(p\) người tạo thêm \(p!\) cách biểu diễn cùng một tập.
Trong bài toán xếp khối (packing problem), biến chiếm chỗ theo từng ô có thể tạo nhiều cách diễn
giải một hình chữ nhật; biến tọa độ và hướng đặt biểu diễn hình học trực tiếp hơn.

Nguyên tắc chung là:
\[
  \text{không biểu diễn thứ tự nếu thứ tự không thuộc nghiệm.}
\]
Sau bước chọn biến, các ràng buộc phá đối xứng xử lý nhóm đối xứng còn lại.
Cách tiếp cận đó giảm cả số biến lẫn số mệnh đề; bước phá đối xứng tiếp tục
giảm số mô hình.

\section{Bốn mẫu phá đối xứng}

\subsection{Cố định một lát chuẩn}

Khi các lớp thời gian hoặc tài nguyên hoán vị tự do, có thể cố định lát đầu.
Với Social Golfer, tuần đầu được chuẩn hóa:
\[
  G_{i,r,1}=1
  \quad\text{khi}\quad
  (r-1)p<i\le rp.
\]
Mọi lịch khả thi có thể đổi tên người chơi và nhóm để thỏa cách sắp này. Những
Các mệnh đề đơn vị đó loại một phần lớn quỹ đạo mà hầu như không tăng chi phí giải
\parencite{kieu2026golfer}.

\subsection{Tiền tố sử dụng và thứ tự ưu tiên giá trị}

Giả sử \(u_b\) cho biết khung chứa \(b\) được dùng. Vì tên các khung chứa có thể hoán vị,
áp đặt
\[
  u_{b+1}\rightarrow u_b.
\]
Các khung chứa được dùng tạo thành một tiền tố. Với các giá trị đối xứng, thứ
tự ưu tiên giá trị yêu cầu lần xuất hiện đầu của giá trị nhỏ hơn đứng trước lần
xuất hiện đầu của giá trị lớn hơn. Hai mẫu này cùng chọn thứ tự theo ``lần sử
dụng đầu'', nhưng tiền tố sử dụng chỉ cần một dãy Boolean nên thường rẻ hơn.

\begin{proposition}[Tiền tố sử dụng không làm mất lớp nghiệm]
\label{prop:usage-prefix}
Giả sử các khung chứa có cùng miền, cùng sức chứa và hàm mục tiêu không phụ
thuộc vào tên khung. Với biến gán \(a_{ib}\) và biến sử dụng \(u_b\), mọi
nghiệm đều có một nghiệm cùng giá trị mục tiêu thỏa
\(u_{b+1}\rightarrow u_b\).
\end{proposition}

\begin{proof}
Cho \(U=\set{b:u_b=1}\) và \(q=\card{U}\). Chọn một song ánh
\(\pi:U\to\set{1,\ldots,q}\), rồi mở rộng \(\pi\) thành một hoán vị trên toàn
bộ tập khung chứa. Định nghĩa
\[
  a'_{i,\pi(b)}=a_{ib},
  \qquad
  u'_{\pi(b)}=u_b .
\]
Mọi biến thuộc bản thân vật \(i\), như tọa độ, phép quay hay kích thước, được
giữ nguyên. Vì sức chứa và các ràng buộc gán được phát biểu đồng nhất cho mọi
khung, đổi tên theo \(\pi\) bảo toàn tính khả thi; vì hàm mục tiêu không đọc
chỉ số khung, giá trị mục tiêu cũng không đổi. Trong nghiệm mới,
\(u'_1=\cdots=u'_q=1\) và \(u'_{q+1}=\cdots=0\), nên toàn bộ implication
\(u'_{b+1}\rightarrow u'_b\) đúng.
\end{proof}

Chứng minh trên cũng chỉ rõ điều kiện không thỏa: nếu khung chứa có sức chứa, chi phí hoặc
miền khác nhau thì hoán vị tùy ý không còn là đối xứng. Khi đó chỉ được áp
dụng tiền tố trong từng lớp khung chứa thật sự đồng nhất.

\subsection{Lex-leader}

Với hai vector đại diện các đối tượng hoán vị được,
\[
  A=\langle a_1,\ldots,a_m\rangle,\qquad
  B=\langle b_1,\ldots,b_m\rangle,
\]
ràng buộc \(A\le_{\mathrm{lex}}B\) chọn thứ tự từ điển. Một phép biểu diễn dùng
các biến bằng nhau trên tiền tố \(e_i\):
\[
  e_i\leftrightarrow e_{i-1}\land(a_i\leftrightarrow b_i),
  \qquad e_0=\top,
\]
và cấm \(e_{i-1}\land a_i\land\neg b_i\). \emph{Lex-leader} (ràng buộc dẫn
đầu theo thứ tự từ điển) mạnh hơn tiền tố sử dụng nhưng tạo thêm một chuỗi biến
phụ cho mỗi cặp được so sánh.

\subsection{Neo hình học và phản xạ}

Với khung chứa hình chữ nhật, các phép phản xạ hoặc quay toàn bố cục có thể bảo
toàn nghiệm. Một vật chuẩn---thường là vật lớn nhất theo diện tích rồi theo
cạnh dài---được neo vào nửa trái hoặc một góc chuẩn:
\[
  2x_a+w_a\le W.
\]
Điều kiện phải khớp với nhóm đối xứng thật. Nếu khung chứa, miền quay hoặc hàm
mục tiêu không bất biến qua phản xạ thì ràng buộc này không hợp lệ.

\begin{keyidea}{Ràng buộc mạnh nhất không mặc nhiên là tốt nhất}
Mỗi ràng buộc phá đối xứng đổi cấu trúc mệnh đề, thứ tự quyết định và các xung
đột học được. Nên so sánh ít nhất ba cấu hình: không phá đối xứng, ràng buộc rẻ
và ràng buộc mạnh. Kết luận dựa trên số bài toán giải được và thời gian, không
chỉ trên số mô hình bị
loại.
\end{keyidea}

\section{Biến quan hệ cho các phép tuyển}

Nhiều bài toán có ràng buộc dạng
\[
  A\lor B\lor C\lor D.
\]
Trong bài toán xếp khối, hai hình \(i,j\) không chồng lấn khi ít nhất một quan hệ đúng:
\[
  L_{ij}\lor L_{ji}\lor B_{ij}\lor B_{ji},
\]
trong đó \(L_{ij}\) nghĩa là \(i\) nằm trái \(j\), còn \(B_{ij}\) nghĩa là
\(i\) nằm dưới \(j\). Trong lập lịch, hai tác vụ dùng cùng tài nguyên cần
\[
  P_{ij}\lor P_{ji},
\]
với \(P_{ij}\) nghĩa là \(i\) kết thúc trước khi \(j\) bắt đầu.

\emph{Witness variable} (biến làm chứng) là biến chỉ ra nhánh nào của phép
tuyển được chọn. Trong mô hình này, biến quan hệ đóng vai trò biến làm chứng
cho một nhánh của phép tuyển. Chúng rút
ràng buộc toàn cục thành các mệnh đề cục bộ quanh cặp đối tượng. Khi một quan
hệ được chọn, lựa chọn này kích hoạt một quan hệ suy diễn kéo theo trên tọa độ hoặc thời gian; khi
miền tọa độ loại quan hệ đó, mệnh đề tuyển buộc bộ giải chọn quan hệ còn lại.

\subsection{Biến quan hệ và giá trị hỗ trợ}

Từ góc nhìn chuyển CSP sang SAT, biến quan hệ nằm giữa biểu diễn trực tiếp và mã
hóa hỗ trợ. Mô hình không chọn toàn bộ giá trị của hai biến; phép biểu diễn chỉ chọn một
\emph{lớp} cặp giá trị,
chẳng hạn ``trái'', ``dưới'' hoặc ``trước''. Mỗi lớp sau đó được liên kết sang
các literal biên của miền. Đây là một phép phân tích thành nhân tử:
\[
  \text{bảng cặp giá trị hợp lệ}
  =
  \bigcup_{q\in Q}\text{lớp hỗ trợ}(q).
\]
Nếu \(|Q|\) nhỏ và mỗi lớp được mô tả bằng vài bất đẳng thức, các biến quan hệ
nén một bảng cấm lớn. Nếu các lớp chồng lấn quá nhiều hoặc không
có mô tả ngắn, biến quan hệ chỉ thêm một tầng trung gian.

\begin{figure}[tb]
\centering
\begin{tikzpicture}[satfig,x=8mm,y=7mm]
  \node[satlabel] at (1.7,3.2) {\(L_{ij}\)};
  \node[satlabel] at (6.7,3.2) {\(L_{ji}\)};
  \node[satlabel] at (1.65,1.3) {\(B_{ij}\)};
  \node[satlabel] at (6.65,1.3) {\(B_{ji}\)};
  \draw[Blue,fill=PaleBlue,thick] (0,2) rectangle (1.5,2.8);
  \node at (.75,2.4) {\(i\)};
  \draw[Teal,fill=PaleTeal,thick] (2,2) rectangle (3.4,2.8);
  \node at (2.7,2.4) {\(j\)};
  \draw[Teal,fill=PaleTeal,thick] (5,2) rectangle (6.4,2.8);
  \node at (5.7,2.4) {\(j\)};
  \draw[Blue,fill=PaleBlue,thick] (6.9,2) rectangle (8.4,2.8);
  \node at (7.65,2.4) {\(i\)};
  \draw[Blue,fill=PaleBlue,thick] (.9,-1) rectangle (2.4,-.2);
  \node at (1.65,-.6) {\(i\)};
  \draw[Teal,fill=PaleTeal,thick] (.9,.1) rectangle (2.4,.9);
  \node at (1.65,.5) {\(j\)};
  \draw[Teal,fill=PaleTeal,thick] (5.9,-1) rectangle (7.4,-.2);
  \node at (6.65,-.6) {\(j\)};
  \draw[Blue,fill=PaleBlue,thick] (5.9,.1) rectangle (7.4,.9);
  \node at (6.65,.5) {\(i\)};
\end{tikzpicture}
\caption{Bốn lớp hỗ trợ của quan hệ không chồng lấn giữa hai hình chữ nhật.
Một literal quan hệ làm chứng cho cả một lớp cặp tọa độ.}
\label{fig:four-spatial-relations}
\end{figure}

\section{Biểu diễn thứ tự và các literal so sánh}

Cho biến nguyên \(z\in\{0,\ldots,U\}\). Biểu diễn thứ tự dùng
\[
  o_{z,t}=1
  \quad\Longleftrightarrow\quad
  z\ge t,
  \qquad t=1,\ldots,U.
\]
Tính đơn điệu được biểu diễn bởi
\[
  o_{z,t+1}\rightarrow o_{z,t}.
\]
Một cận \(z\ge a\) là mệnh đề đơn vị \(o_{z,a}\); cận \(z\le b\) là
\(\neg o_{z,b+1}\). Phép so sánh dịch chuyển
\[
  x_i+w_i\le x_j
\]
trở thành các mệnh đề cấm mọi cặp ngưỡng mâu thuẫn. Khi \(w_i\) là hằng, cấu
trúc dịch chuyển này có tính địa phương cao và phù hợp với bài toán xếp khối (packing problem)
\parencite{soh2010binpacking,kieu2025strip,kieu2025bin}.

Với lập lịch, cùng kỹ thuật biểu diễn
\[
  s_i+p_i\le s_j.
\]
Nhờ vậy, biến quan hệ \(P_{ij}\) có thể liên kết trực tiếp với các literal thứ
tự của \(s_i,s_j\), thay vì tạo một mạch cộng tổng quát
\parencite{kieu2025scheduling}.

\begin{workedexample}{Liên kết biểu diễn trực tiếp--thứ tự cho một giá trị nguyên}
Cho \(z\in\{0,1,2,3,4\}\). Biểu diễn trực tiếp có \(d_a\Leftrightarrow(z=a)\);
biểu diễn thứ tự có \(o_t\Leftrightarrow(z\ge t)\). Hai góc nhìn được nối bởi
\[
  d_a\leftrightarrow o_a\land\neg o_{a+1},
  \qquad o_0\equiv\top,\quad o_5\equiv\bot.
\]
Nếu bộ giải suy ra \(o_3=1,o_4=0\), các mệnh đề liên kết buộc \(d_3=1\), rồi
ràng buộc đúng một trên \(d_0,\ldots,d_4\) loại mọi giá trị khác. Theo chiều
ngược, \(d_2=1\) buộc \(o_1=o_2=1\) và \(o_3=o_4=0\), khiến các ràng buộc cận
đọc được ngay \(z=2\). Ví dụ này cho thấy lợi ích của liên kết hai chiều không
nằm ở số mô hình mà ở đường truyền suy luận giữa hai giao diện.
\end{workedexample}

\section{Liên kết một chiều hay hai chiều}

Giả sử \(q\) là biến quan hệ và \(\Phi\) là công thức miền gốc biểu diễn quan
hệ đó. Có ba mức liên kết:
\[
  q\rightarrow\Phi,\qquad
  \Phi\rightarrow q,\qquad
  q\leftrightarrow\Phi.
\]
Chiều đầu đủ cho tính đúng nếu \(q\) chỉ được dùng làm biến làm chứng được chọn:
bộ giải không thể đặt \(q=1\) khi quan hệ hình học sai. Chiều ngược giúp miền
tọa độ suy ra \(q\); liên kết hai chiều cho phép lan truyền qua cả hai biểu
diễn.

Lựa chọn phụ thuộc vào cách mô-đun sử dụng đọc biến:

\begin{itemize}
  \item Nếu mệnh đề \(q_1\lor\cdots\lor q_r\) yêu cầu một biến làm chứng và không
  nơi nào đọc \(\neg q\), quan hệ kéo theo tiến thường đủ.
  \item Nếu \(\neg q\) được dùng để kết luận quan hệ đối nghịch, phải có
  đặc tả mạnh hơn.
  \item Nếu nhiều ràng buộc và hàm mục tiêu chia sẻ \(q\), quan hệ tương
  đương giúp
  tránh các trạng thái phụ không mang nghĩa.
\end{itemize}

\begin{example}
Với hai tác vụ cùng máy, đặt \(P_{ij}\lor P_{ji}\). Hai mệnh đề liên kết
\[
  P_{ij}\rightarrow(s_i+p_i\le s_j),
  \qquad
  P_{ji}\rightarrow(s_j+p_j\le s_i)
\]
đã bảo đảm chúng không chồng lấn. Nếu quan hệ trước--sau từ một ràng buộc khác cho
biết \(s_i+p_i\le s_j\), chiều ngược
\((s_i+p_i\le s_j)\rightarrow P_{ij}\) có thể truyền quan hệ này sang các
mô-đun khác.
\end{example}

\subsection{Lan truyền qua vòng liên kết}

Phép biểu diễn kết hợp tạo một vòng suy luận:
\[
  \text{lựa chọn chính xác}
  \rightarrow
  \text{các ngưỡng}
  \rightarrow
  \text{quan hệ}
  \rightarrow
  \text{ngưỡng của biến khác}.
\]
Mỗi chiều liên kết đóng một đường suy luận trong vòng trên. Phép đánh giá bắt
đầu bằng danh sách các suy luận miền mong muốn, sau đó tìm tập chiều tối thiểu
tạo được chúng. Đặc tả này có thể kiểm bằng các phép gán một phần và lan truyền
đơn vị như trong Chương 2.

Việc giữ nhiều góc nhìn cũng liên quan đến biểu diễn lại có cấu trúc: biến phụ có
giá trị khi đặt tên cho một lớp hỗ trợ được nhiều ràng buộc dùng lại, thay
vì chỉ thay một biểu thức xuất hiện một lần
\parencite{haberlandt2023aux}.

\begin{figure}[tb]
\centering
\begin{tikzpicture}[satfig,node distance=17mm]
  \node[satbox] (d) {lựa chọn trực tiếp\\\(d_a\)};
  \node[satbox,right=of d] (o) {literal thứ tự\\\(o_t\)};
  \node[satbox,right=of o] (q) {quan hệ\\\(q\)};
  \node[satbox,below=of o] (v) {miền của\\biến khác};
  \draw[satedge] (d)--node[above,satlabel]{giá trị \(\to\) cận}(o);
  \draw[satedge] (o)--node[above,satlabel]{hỗ trợ}(q);
  \draw[satedge] (q)--node[right,satlabel]{dịch chuyển}(v);
  \draw[satedge] (v)--node[left,satlabel]{loại cận}(d);
  \draw[satdash,bend left=20] (o) to node[below,satlabel]{chiều ngược}(d);
\end{tikzpicture}
\caption{Vòng liên kết trong phép biểu diễn kết hợp. Thiếu một chiều có thể làm
đứt một đường lan truyền đơn vị dù công thức vẫn đúng theo phép chiếu.}
\label{fig:channeling-loop}
\end{figure}

\section{Phép quay và kích thước có điều kiện}

Khi vật \(i\) được phép quay, biến \(\rho_i\) chọn hướng đặt:
\[
  (w_i(\rho_i),h_i(\rho_i))
  =
  \begin{cases}
    (w_i,h_i),&\rho_i=0,\\
    (h_i,w_i),&\rho_i=1.
  \end{cases}
\]
Các literal quan hệ được chia sẻ cho hai hướng đặt và được liên kết với phép
so sánh dịch chuyển dưới điều kiện gác \(\rho_i,\rho_j\). Ví dụ, \(L_{ij}\)
sinh bốn họ mệnh đề, mỗi họ ứng với một tổ hợp hướng đặt còn hợp lệ. Các tổ
hợp không vừa khung chứa bị loại bằng mệnh đề đơn vị hoặc mệnh đề nhị phân
trước khi sinh quan hệ cặp.

Cách tổ chức này tách ba tầng:
\[
  \text{hướng đặt}
  \longrightarrow
  \text{kích thước hiệu dụng}
  \longrightarrow
  \text{quan hệ không chồng lấn}.
\]
Cấu trúc này cũng cho phép tái sử dụng cùng biến \(\rho_i\) trong cận miền, quy tắc trội
và bộ giải mã.

\section{Đồ thị xung đột và tính địa phương}

Đồ thị xung đột \(H=(V,E)\) chọn chính xác các cặp cần biến quan hệ:
\(\{i,j\}\in E\) khi hai đối tượng có thể cạnh tranh cùng tài nguyên, chồng
lấn hình học hoặc vi phạm một cận cục bộ.

Trong lập lịch, cạnh xuất hiện khi hai tác vụ có thể cùng dùng một tài nguyên
và miền thời gian của chúng giao nhau. Trong bài toán xếp khối, mọi cặp vật cùng khung
chứa là ứng viên, nhưng cận miền và quan hệ trội có thể cố định trước một số
quan hệ. Trong phân bổ tần số, cạnh là cặp máy phát có yêu cầu khoảng cách nhãn.

Thiết kế theo đồ thị đem lại hai lợi ích: kích thước tỉ lệ với số tương tác
thật, và mệnh đề học được thường nằm quanh một lân cận nhỏ. Đây là một tiêu chí
cụ thể về mức độ phù hợp với bộ giải, thay vì chỉ nói phép biểu diễn ``gọn''.

\section{Kiến trúc kết hợp}

Một mô hình SAT lớn nên ghép các mô-đun theo thứ tự sau:

\begin{enumerate}
  \item chọn biến chính loại các chỉ số không mang nghĩa;
  \item tạo miền thứ tự/trực tiếp cho thời gian, tọa độ hoặc nhãn;
  \item tạo biến quan hệ trên đồ thị xung đột;
  \item liên kết mỗi quan hệ với các literal miền;
  \item thêm ràng buộc đếm hoặc bộ đếm chia sẻ cho tài nguyên;
  \item thêm các ràng buộc phá đối xứng có chứng minh quỹ đạo;
  \item gắn hàm mục tiêu bằng các giả định hoặc mệnh đề mềm.
\end{enumerate}

Các nghiên cứu về lập lịch, xếp khối dải/khung chứa và Social Golfer trong
\textcite{kieu2025thesis} cho thấy ba lựa chọn đầu---biến chính, quan hệ và
đối xứng---quyết định cấu trúc CNF trước khi chọn bộ giải. Các chương sau dùng
chính kiến trúc này cho bốn họ bài toán.

\section{Bài học thiết kế}

\begin{summarybox}
\begin{enumerate}
  \item Loại đối xứng biểu diễn trước khi thêm các mệnh đề phá đối xứng.
  \item Một breaker hợp lệ phải giữ ít nhất một đại diện của mọi quỹ đạo.
  \item Biến quan hệ là biến làm chứng có ngữ nghĩa, giúp cục bộ hóa phép tuyển.
  \item Liên kết mạnh đến đâu phụ thuộc vào cách biến quan hệ được đọc.
  \item Biểu diễn thứ tự làm các cận dịch chuyển trở thành mệnh đề đơn giản.
  \item Đồ thị xung đột xác định nơi thật sự cần biến phụ.
\end{enumerate}
\end{summarybox}

\subsection*{Câu hỏi nghiên cứu}
\begin{enumerate}
  \item Chứng minh tiền tố sử dụng giữ một đại diện khi các khung chứa hoàn
  toàn giống nhau.
  \item Viết liên kết một chiều cho bốn quan hệ không chồng lấn của hai hình.
  \item Cho một ràng buộc phản xạ trong bài toán xếp khối (packing problem) dải và nêu điều kiện
  để ràng buộc đó hợp lệ.
\end{enumerate}
